\documentclass[11pt,reqno]{article}
\usepackage{amsmath,amsthm,amssymb}
\usepackage[margin=1.1in]{geometry}
\usepackage{booktabs}
\emergencystretch=3em
\newtheorem{theorem}{Theorem}
\newtheorem{lemma}[theorem]{Lemma}
\newtheorem{corollary}[theorem]{Corollary}
\newtheorem{proposition}[theorem]{Proposition}
\newtheorem{conjecture}[theorem]{Conjecture}
\theoremstyle{definition}
\newtheorem{definition}[theorem]{Definition}
\newtheorem{example}[theorem]{Example}
\newtheorem{remark}[theorem]{Remark}
\newcommand{\ggbrack}[2]{\genfrac{[}{]}{0pt}{}{#1}{#2}_q}
\title{\bf A shift law for central gaps of Gaussian binomial coefficients,\\ with sharp unimodality thresholds for a two-rectangle subtraction family}
\author{Gerald W.\ Hu}
\date{November 13, 2026}
\begin{document}
\maketitle

\begin{abstract}
Let $a_l(n)$ denote the number of partitions of $n$ fitting in the $l\times 6$ rectangle, i.e.\ the coefficient of $q^n$ in the Gaussian polynomial $\ggbrack{l+6}{6}$, and let $g_6(l)=a_l(3l)-a_l(3l-1)$ be the central gap of $\ggbrack{l+6}{6}$ (the symmetry center is $3l$). We prove the \emph{Shift Law}:
\[
g_6(l)=g_6(l-15)\qquad\text{for every odd } l\ge 17,
\]
so the anomalously small odd-parity central gaps are exactly the even-parity gaps transported by the constant shift $15=\binom{6}{2}$. The proof is a complete symbolic argument: an inclusion--exclusion identity expresses $g_6(l)$, for $l\ge 12$, as an explicit integer combination of the $\{2,3,4,5,6\}$-denumerant $Q_6$; $Q_6$ is certified to be a quasipolynomial of period $60$ valid for all $n\ge0$; and the Shift Law then reduces to $30$ explicit polynomial identities between residue-class branch polynomials, each verified exactly. Combining the Shift Law with a certified central-dominance equivalence, we obtain exact sharp unimodality classifications for the subtraction family $F^{(6,b)}_l(q)=\ggbrack{l+6}{6}-q^{3l-b^2/8}\ggbrack{b}{b/2}$ at $b\in\{8,12,16\}$: $F^{(6,b)}_l$ is unimodal if and only if $g_6(l)\ge\Delta_b$ (with $\Delta_8=1$, $\Delta_{12}=3$, $\Delta_{16}=7$), and the exceptional sets are exactly $\{3,5,7,9,11,13\}$, $\{7,9,\dots,19\}$, and $\{11,13,\dots,25\}$ (odd $l$), with sharp thresholds $l_0=14,20,26$. The equivalence is false in general: we exhibit the certified counterexample $(k,b,l)=(6,24,46)$, where the central inequality passes with equality and an interior adjacent difference fails. The constant $15$ is the degree of the odd generator of the classical invariant ring of binary sextics, whose Hilbert series makes the shift immediate; what remains open is a partition-level bijective explanation, and we document precisely why the standard staircase, complementation, and rim-hook transports do not supply one. All infinite computer-assisted content is reduced to explicit finite certificates in exact integer arithmetic; the accompanying scripts reproduce every decisive identity, and deliberate certificate corruption is detected.
\end{abstract}

\section{Introduction}\label{sec:intro}
The Gaussian polynomial (or $q$-binomial coefficient) $\ggbrack{n}{k}$ enumerates partitions fitting in a $k\times(n-k)$ rectangle. Its unimodality, proved by Sylvester in 1878 \cite{sylvester}, remains a central object of enumerative combinatorics; see \cite{stanley} for the standard reference and \cite{ohara} for a constructive proof.

Two lines of recent work motivate this paper. First, Stanley and Zanello \cite{SZ} conjectured unimodality (with nonnegativity) for differences
\[
\ggbrack{m}{k}-q^{\,k(m-b)/2+b-2k+2}\ggbrack{b}{k-2}
\]
in a precisely specified range, proving the cases $k\le 5$; Koutschan, Uncu and Wong \cite{KUW} settled the cases $k=6,7$ of that conjecture by quantifier elimination for quasipolynomials (cylindrical algebraic decomposition). Second, Reiner and Stanton \cite{RS} proved, among other things, a two-step propagation theorem: for fixed even $k$ and a fixed central subtractand, unimodality of the difference for two consecutive base values of the rectangle width propagates along the corresponding parity class. Their Theorem 5 in particular implies the special case $b=k+2$ of the Stanley--Zanello family. Exact thresholds and exact exceptional sets, however, are not accessible from either the propagation mechanism or the eventual-in-$m$ conjecture: both are silent about which parity class is slow and where it crosses.

This paper fixes the box height $k=6$ and studies the \emph{central gap}
\[
g_6(l)\;=\;a_l(3l)-a_l(3l-1),
\]
the last ascent of $\ggbrack{l+6}{6}$ at its symmetry center $3l$. The function $g_6$ controls unimodality of the centered subtraction family
\[
F^{(6,b)}_l(q)\;=\;\ggbrack{l+6}{6}-q^{3l-b^2/8}\ggbrack{b}{b/2},\qquad b\equiv 0\pmod 4,
\]
whose two summands are both symmetric about degree $3l$ (see \S\ref{sec:setup}). Three facts, each independently certified, organize the story:
\begin{enumerate}
\item \textbf{Parity obstruction.} For odd $l\le 13$ one has $g_6(l)=0$, while $g_6$ grows without bound on even $l$ (degree-$3$ branch polynomials, Appendix~\ref{app:cert}): the odd residue class is anomalously slow. (This is the $k=6$ instance of a slow-class pattern $kl\equiv 2\pmod 4$ that is \emph{observed}, not proved, for $k>6$; see \S\ref{sec:why15}.)
\item \textbf{Shift Law} (Theorem~\ref{thm:shift}): for every odd $l\ge17$, $g_6(l)=g_6(l-15)$. The slow class is \emph{not} asymptotically smaller; it is the fast class, translated by $15=\binom 62$.
\item \textbf{Central dominance over the certified range} (Theorem~\ref{thm:equivalence}): for $b\in\{8,12,16\}$ and every admissible $l$, $F^{(6,b)}_l$ is unimodal if and only if $g_6(l)\ge\Delta_b$, where $\Delta_b$ is the last ascent of $\ggbrack{b}{b/2}$.
\end{enumerate}
Together these give complete, exact, and sharp unimodality classifications (Theorem~\ref{thm:classification}): the exceptional $l$ are finite, contiguous odd segments, and the thresholds are
\[
l_0(6,8)=14,\qquad l_0(6,12)=20,\qquad l_0(6,16)=26 .
\]
The equivalence in (3) is \emph{not} a theorem of universal scope: Proposition~\ref{prop:ce} exhibits the certified counterexample $(6,24,46)$, where the central inequality holds with equality but an interior adjacent difference fails. This is why the classification is stated with an explicit domain.

\subsection*{Method and the status of machine verification}
The proof technique is quasipolynomial: the coefficients of $\ggbrack{l+6}{6}$ are, after inclusion--exclusion, values of the denumerant $Q_6(n)$ (partitions of $n$ into parts in $\{2,3,4,5,6\}$), which is a quasipolynomial of period $60$ and degree $4$ \cite{KR}. The two infinite steps in the proof are therefore certified as follows.
\begin{itemize}
\item $Q_6$ equals its period-$60$ branch interpolants for \emph{all} $n\ge0$: verified by a constant-coefficient recurrence check (symbolic, $60/60$ residue classes) plus a finite induction base $0\le n\le119$ (Lemma~\ref{lem:qp-cert}).
\item The Shift Law reduces to $30$ explicit polynomial identities $G_r(t)\equiv G_{r-15}(t)$ between branch polynomials with exact rational coefficients; all $30$ verify symbolically (Theorem~\ref{thm:shift}'s proof, Appendix~\ref{app:cert} for sample data).
\item The central-dominance equivalence uses $60h$ tail certificates per family ($h=b^2/8$) built on a sound minimum test for integer polynomials of degree at most $3$ (Lemma~\ref{lem:cubic}), plus exact finite computation for $l\le 120$.
\end{itemize}
Following the epistemology of \cite{KUW}, we label in \S\ref{sec:status} exactly which statements are hand proofs, which are machine-verified certificates over explicit finite data, and which are numerically supported but unproved. No numerical observation is promoted to a theorem anywhere in this paper.

\subsection*{Relation to prior work}
The quasipolynomial method with a positivity certificate is standard technology \cite{KUW,KKM}; the reduction of unimodality of $F^{(k,b)}_l$ to finitely many fixed adjacent-offset inequalities is implicit in \cite{RS} and explicit in the companion papers \cite{M06,M07}. The Shift Law is a consequence of the classical Hilbert series of the invariant ring of binary sextics and follows from Appendix A.2 of \cite{KKM} (identified in external review, 2026-11-13; \S\ref{sec:related}); our residue-class enumeration gives a self-contained elementary proof of the same identity. The contributions of this paper are the central-gap equivalence over an explicit certified domain, the exact exceptional-set classifications with sharp thresholds, and the certified failure at $(6,24)$; we make no stronger novelty claim for these than ``not found in the sources cited.'' A partition-level explanation of the constant $15$ is open (\S\ref{sec:why15}).

\section{Setup, definitions, and statements}\label{sec:setup}
Throughout, $l,k,b$ are integers. For $n\ge k\ge 0$,
\[
\ggbrack{n}{k}=\frac{(1-q)(1-q^2)\cdots(1-q^{n})}{(1-q)\cdots(1-q^{k})\,(1-q)\cdots(1-q^{n-k})}
\]
is a polynomial in $q$ with nonnegative integer coefficients, palindromic of degree $k(n-k)$; it is the generating function of partitions fitting in the $k\times(n-k)$ rectangle. We write
\[
a_l(n)\;=\;[q^n]\,\ggbrack{l+6}{6}\;=\;\#\{\lambda\vdash n:\ \ell(\lambda)\le 6,\ \lambda_1\le l\},
\]
with $a_l(n)=0$ for $n<0$ or $n>6l$. The polynomial $\ggbrack{l+6}{6}$ has degree $6l$ and is symmetric about degree $3l$, so the \emph{central gap}
\[
g_6(l)\;:=\;a_l(3l)-a_l(3l-1)
\]
is the last first-difference before the center; it is an integer $\ge 0$ by Sylvester's unimodality theorem.

Fix $b\equiv 0\pmod 4$ and put $h=b^2/8$, and for $3l\ge h$ (the \emph{admissible} range) define
\[
F^{(6,b)}_l(q)\;=\;\ggbrack{l+6}{6}-q^{s}\,\ggbrack{b}{b/2},\qquad s=3l-h .
\]
Both summands are symmetric about degree $3l$: the second because $q^s\ggbrack{b}{b/2}$ has degree $s+2h=3l+h$ and is palindromic about $s+h=3l$. Write $b_j=[q^j]\ggbrack{b}{b/2}$ and set $b_j=0$ for $j\notin[0,2h]$. The \emph{terminal ascent} of the subtractand is
\begin{equation}\label{eq:Delta}
\Delta_b\;:=\;b_h-b_{h-1}
\;=\;[q^h]\ggbrack{b}{b/2}-[q^{h-1}]\ggbrack{b}{b/2} .
\end{equation}
Exact values, computed from the definition by the scripts of \S\ref{sec:status}:
\begin{equation}\label{eq:Deltatable}
\Delta_8=1,\qquad \Delta_{12}=3,\qquad \Delta_{16}=7,\qquad \Delta_{20}=24,\qquad \Delta_{24}=127 .
\end{equation}

\begin{theorem}[Shift Law]\label{thm:shift}
For every odd integer $l\ge 17$,
\[
g_6(l)\;=\;g_6(l-15).
\]
The statement is one-directional. For even $l\ge2$ the theorem applied to the odd parameter $l+15$ gives $g_6(l+15)=g_6(l)$ (e.g.\ $g_6(24)=g_6(39)=29$), whereas the reverse comparison $g_6(l)=g_6(l-15)$ fails in general on even $l$ ($g_6(24)=29\ne 0=g_6(9)$); there is no period-$15$ symmetry.
\end{theorem}

\begin{theorem}[Central-gap equivalence]\label{thm:equivalence}
Let $(k,b)\in\{(6,8),(6,12),(6,16)\}$. For every admissible $l$ (i.e.\ $3l\ge b^2/8$), $F^{(6,b)}_l(q)$ is unimodal if and only if
\[
g_6(l)\;\ge\;\Delta_b .
\]
\end{theorem}

\begin{theorem}[Exact classification]\label{thm:classification}
For admissible $l\ge 1$, the set of $l$ for which $F^{(6,b)}_l$ is not unimodal is exactly
\[
E(6,8)=\{3,5,7,9,11,13\},\qquad E(6,12)=\{7,9,11,13,15,17,19\},
\]
\[
E(6,16)=\{11,13,15,17,19,21,23,25\}.
\]
In particular the thresholds $l_0(6,b):=\min\{l^\ast:\ F^{(6,b)}_l\ \text{is unimodal for all admissible } l\ge l^\ast\}$ are $l_0(6,8)=14$, $l_0(6,12)=20$, $l_0(6,16)=26$, and each is sharp. With the convention $g_6(0):=1$ and
\[
l_{0e}(6,b):=\min\{\,l\ge 0\ \text{even}:\ g_6(l)\ge\Delta_b\,\},
\]
the thresholds obey $l_0(6,b)=l_{0e}(6,b)+14$, i.e.\ the last exceptional $l$ is $l_{0e}+13=l_{0e}-2+15$, the translate of the largest even parameter $l_{0e}-2$ with $g_6<\Delta_b$ (whenever $l_{0e}\ge 2$; for $b=8$, $l_{0e}=0$ and the last exception $l=13$ is verified directly).
\end{theorem}

\begin{proposition}[Failure of universal central dominance]\label{prop:ce}
Let $h=72$, $s=3l-72$, and $F^{(6,24)}_l(q)=\ggbrack{l+6}{6}-q^{s}\ggbrack{24}{12}$. For $l=46$ (admissible: $s=66\ge0$) the central inequality holds with equality,
\[
g_6(46)-\Delta_{24}=0,
\]
yet $F^{(6,24)}_{46}$ is not unimodal: its adjacent difference at offset $j=h-6$ equals $-86$, i.e.\
\[
[q^{133}]F^{(6,24)}_{46}-[q^{132}]F^{(6,24)}_{46}=-86 .
\]
Equally, $l=48$ has $g_6(48)-\Delta_{24}=14>0$ and an adjacent difference of $-11$ at offset $j=h-2$. Hence no statement of the form ``$F^{(6,b)}_l$ unimodal $\Leftrightarrow g_6(l)\ge\Delta_b$'' can hold for all $b$.
\end{proposition}

Proposition~\ref{prop:ce} is a single exact integer computation; the accompanying script \texttt{counter\allowbreak example\_6\_24.py} reproduces both lines from the definition of the Gaussian polynomial. The mechanism behind it is analyzed in \S\ref{sec:cubic}: central dominance holds exactly when the ascent profile of the large rectangle outruns that of the subtractand at every interior offset, and for wide subtractands ($b=24$ at $k=6$) it marginally fails.

\begin{remark}[Small parameters and conventions]\label{rem:conventions}
(i) Admissibility $3l\ge h$ is exactly the condition that $F^{(6,b)}_l$ be an honest polynomial rather than a Laurent polynomial; below it the window reduction of Lemma~\ref{lem:window} does not apply and we make no claims.
(ii) The degenerate identity case $k=l=b/2$ (here $l=6$, $b=12$) gives $s=0$ and $F^{(6,12)}_6\equiv 0$, the zero polynomial, which is unimodal; the equivalence of Theorem~\ref{thm:equivalence} correctly records this since $g_6(6)=3=\Delta_{12}$.
(iii) The convention $g_6(0)=1$ (from $a_0(0)=1$, $a_0(-1)=0$) makes the Shift Law instance $l=15$ true ($g_6(15)=1=g_6(0)$), but $l=15$ and $l=0$ are kept outside Theorem~\ref{thm:shift} because the box degenerates.
\end{remark}

\section{The window reduction}\label{sec:window}
All three theorems rest on reducing unimodality to finitely many adjacent-offset inequalities. This step is a hand proof using no computation.

For $j\in\{-1,0,\dots,h-1\}$ define the \emph{margin}
\[
M_j(l)\;:=\;a_l(s+j+1)-a_l(s+j)\;-\;\bigl(b_{j+1}-b_j\bigr),
\qquad s=3l-h .
\]
Note $M_{h-1}(l)=g_6(l)-\Delta_b$, since $s+h=3l$.

\begin{lemma}[Window reduction]\label{lem:window}
Let $b\equiv 0\pmod 4$ and let $l$ be admissible. Then $F^{(6,b)}_l$ is unimodal if and only if
$M_j(l)\ge 0$ for every $j\in\{-1,\dots,h-1\}$.
\end{lemma}
\begin{proof}
Let $F[n]:=[q^n]F^{(6,b)}_l$ and $D(n):=F[n+1]-F[n]$. Both summands of $F$ are symmetric about $3l$ and $F$ has degree $6l$, so $F[n]=F[6l-n]$ and $D(3l+u)=-D(3l-u-1)$; it suffices to control $D(n)$ for $0\le n\le 3l-1$ and to establish nonnegativity of all coefficients.

\emph{Case $n\le s-2$:} then $n+1-s\le -1$, so the $b$-term contributes $b_{n+1-s}-b_{n-s}=0-0=0$, and $D(n)=a_l(n+1)-a_l(n)\ge0$ by Sylvester's unimodality of $\ggbrack{l+6}{6}$, since $n+1\le s-1<3l$.

\emph{Case $s-1\le n\le 3l-1$:} write $j=n-s\in[-1,\,h-1]$; then $D(n)=M_j(l)$ by definition.

\emph{Nonnegativity:} for $n<s$, $F[n]=a_l(n)\ge0$; for $n>s+2h$, $F[n]=a_l(n)\ge0$; and for $s\le n\le 3l$, telescoping from $n=s$ gives $F[n]\ge F[s]=a_l(s)-b_0=a_l(s)-1\ge0$, because $0\le s\le 6l$ implies $a_l(s)\ge1$ (write $s=6q+r$, $0\le r<6$; the parts $q{+}1,\dots,q{+}1,q,\dots,q$ with $r$ copies of $q{+}1$ fit the $l\times6$ box and sum to $s$; admissibility gives $q\le l$, and $r>0$ forces $q\le l-1$). Symmetry covers $n>3l$. Hence coefficientwise nonnegativity holds whenever all $M_j\ge0$.

Finally, if all $M_j\ge0$, then $D(n)\ge0$ for all $n\le 3l-1$, so $F$ increases up to degree $3l$ and, by palindromy, decreases after it: $F$ is unimodal. Conversely, suppose $F$ is unimodal. Its maximal mode interval $[m_1,m_2]$ (the set of maximizers, contiguous by unimodality) is mapped to itself under $n\mapsto 6l-n$, since $F[n]=F[6l-n]$; hence $m_1+m_2=6l$ and $m_1\le 3l\le m_2$, so $F$ is nondecreasing on $[0,3l]$: a negative $D(n)$ with $n\le 3l-1$ is impossible. In particular some $M_j<0$ (a negative $D$ at $n=s+j\le 3l-1$) contradicts unimodality.
\end{proof}

\begin{remark}\label{rem:ce-mech}
The converse in the proof shows exactly what fails in Proposition~\ref{prop:ce}: the mode of $F^{(6,24)}_{46}$ drifts off-center while the central difference is exactly $0$.
\end{remark}

\section{The denumerant certificate}\label{sec:qp}
Let $Q_6(n)$ be the number of partitions of $n$ into parts belonging to $\{2,3,4,5,6\}$, with $Q_6(n)=0$ for $n<0$ and $Q_6(0)=1$; its generating function is
\[
\sum_{n\ge0}Q_6(n)q^n=\prod_{i=2}^{6}\frac1{1-q^i}.
\]

\begin{lemma}[Quasipolynomial structure, classical]\label{lem:qp-classical}
$Q_6$ is a quasipolynomial in $n$ of degree $4$ and period dividing $60=\mathrm{lcm}(2,3,4,5,6)$, valid for every $n\ge0$.
\end{lemma}
\begin{proof}
Partial fractions of $\prod_{i=2}^6(1-q^i)^{-1}$: all poles are roots of unity of order dividing $60$, and extracting coefficients of $c(1-q/\zeta)^{-r}$ gives $\binom{n+r-1}{r-1}\zeta^{-n}$, a polynomial in $n$ of degree $r-1\le 4$ times a $60$-periodic factor. There is no polynomial part, so the identity holds from $n=0$ onward. (This is the standard denumerant theorem; cf.\ \cite{KR}.)
\end{proof}

Lemma~\ref{lem:qp-classical} is structural: it does not produce the branch polynomials. For the certificate we use the following effective equivalent, which is what the scripts verify.

Let $F_r(t)$, $r\in\{0,\dots,59\}$, be the unique polynomial of degree $\le4$ interpolating $Q_6$ at $n=r+60t$, $t=0,1,2,3,4$, and set $F(n):=F_{n\bmod 60}(\lfloor n/60\rfloor)$. Expanding $\prod_{i=2}^6(1-q^i)=\sum_d D_d q^d$ (depth $20$, integer coefficients) gives the recurrence $\sum_d D_d\,Q_6(m-d)=0$ for $m\ge1$.

\begin{lemma}[Effective certificate]\label{lem:qp-cert}
Suppose
\begin{enumerate}
\item[(i)] for every residue $r\bmod 60$, the substitution $n=r+60t$ turns the recurrence into an identity of polynomials in $t$:
$\sum_d D_d\,F_{(r-d)\bmod 60}\bigl(t+\lfloor (r-d)/60\rfloor\bigr)\equiv 0$;
\item[(ii)] $F(n)=Q_6(n)$ for all $0\le n\le119$.
\end{enumerate}
Then $Q_6(n)=F(n)$ for all $n\ge0$.
\end{lemma}
\begin{proof}
Induction on $n$. Write $n=r+60t$. In the recurrence $\sum_d D_d Q_6(n-d)=0$ the coefficient of $Q_6(n)$ is $D_0=1$, and every other index $n-d$ satisfies $n-20\le n-d<n$. If $n\ge 120$, then $n-d\ge 100\ge0$ and the induction hypothesis applies to all $Q_6(n-d)$; the same relation with $Q_6$ replaced by $F$ holds at $(r,t)$ by (i); hence $Q_6(n)=F(n)$. The base $n<120$ is (ii).
\end{proof}

\begin{remark}\label{rem:S1}
The script \texttt{verify\_shift.py} (step S1) checks (i) symbolically in all $60$ classes (each check is a polynomial identity with exact rational coefficients, verified by expansion) and (ii) by exact dynamic programming for $0\le n\le119$. This is a finite amount of explicit data; see Appendix~\ref{app:cert}.
\end{remark}

\section{The Shift Law}\label{sec:shift}

\subsection{Reduction of $g_6$ to $Q_6$}
\begin{lemma}[Inclusion--exclusion identity]\label{lem:star}
For every integer $l\ge 12$,
\begin{equation}\label{eq:star}
g_6(l)\;=\;Q_6(3l)\;-\;\sum_{i=1}^{6}Q_6(2l-i)\;+\;\sum_{1\le i<j\le 6}Q_6(l-i-j) .
\tag{$*$}
\end{equation}
\end{lemma}
\begin{proof}
From $\ggbrack{l+6}{6}=\prod_{i=1}^6(1-q^{l+i})\big/\prod_{i=1}^6(1-q^i)$, expanding the numerator and writing $p(m)=p_{\le 6}(m)$ for the number of partitions of $m$ into parts in $\{1,\dots,6\}$ (and $0$ for $m<0$) gives, for every $n$,
\[
a_l(n)=\sum_{S\subseteq\{1,\dots,6\}}(-1)^{|S|}\,p\bigl(n-|S|\,l-\sigma(S)\bigr),
\qquad \sigma(S)=\sum_{i\in S}i .
\]
Apply this at $n=3l$ and $n=3l-1$ and subtract. For $|S|\ge3$ the arguments are $(3-|S|)l-\sigma(S)$ and $(3-|S|)l-\sigma(S)-1$, both negative (at $|S|=3$ the largest is $-6$ for $l=12$, since $\sigma(S)\ge6$), so all subsets of size $\ge3$ contribute to neither difference. For $|S|=2$ the arguments are $l-i-j$ and $l-i-j-1$; and $|S|=1$ gives $2l-i$, $2l-i-1$; all are $\ge 0$ for $l\ge 12$ (at $l=12$ the extremal pair has $i+j=11$, where the two arguments are $1$ and $0$). Subtracting the two identities and using $Q_6(m)=p(m)-p(m-1)$ (which follows from $(1-q)\prod_{i=1}^6(1-q^i)^{-1}=\prod_{i=2}^6(1-q^i)^{-1}$) yields \eqref{eq:star}.
\end{proof}

\begin{remark}
The threshold $l\ge12$ is the point at which the \emph{survival pattern} of subsets used above is stable; identity \eqref{eq:star} happens to hold already for all $l\ge1$ (audited by exact DP to $l=80$), but only $l\ge12$ is needed and used.
\end{remark}

\subsection{Branch polynomials and the proof of Theorem~\ref{thm:shift}}
By Lemma~\ref{lem:star}, for $l\ge12$ the value $g_6(l)$ is an explicit $\{-1,0,+1\}$-linear combination of values of $Q_6$ at arguments $3l$, $2l-i$, $l-i-j$. Fix an odd residue $r\bmod 60$ and write $l=60t+r$. Each argument $\alpha l+\beta$ (with $\alpha\in\{1,2,3\}$, $\beta\in\mathbb Z$) maps to the branch
\[
Q_6(\alpha l+\beta)\;=\;F_{(\alpha r+\beta)\bmod 60}\Bigl(\alpha\,t+\Bigl\lfloor\frac{\alpha r+\beta}{60}\Bigr\rfloor\Bigr),
\]
exact integer floor bookkeeping, valid whenever the argument is nonnegative. Substituting and summing produces, for each odd $r$, an explicit polynomial $G_r(t)\in\mathbb Q[t]$ of degree $\le 4$ such that
\[
g_6(60t+r)\;=\;G_r(t)\qquad\text{for all } t \text{ with } 60t+r\ge 27 .
\]
(The offset $\lfloor(\alpha r+\beta)/60\rfloor$ can shift the branch quotient by $-1$; the certificate evaluates at $t$ large enough that all arguments are nonnegative, and for the wrapped classes $r<15$ the comparison class is $r+45$ at quotient $t-1$.)

\begin{proof}[Proof of Theorem~\ref{thm:shift}]
\emph{Step 1 (branch identities).} For each of the $30$ odd residues $r\bmod 60$, the identity
\begin{equation}\label{eq:branch}
G_r(t)\;\equiv\;G_{r-15}(t)
\end{equation}
holds as a polynomial identity in $t$ with exact rational coefficients. This is the content of step S3 of \texttt{verify\_shift.py}: $30$ symbolic expansions, all zero (Appendix~\ref{app:cert} lists the data, e.g.\ $G_{27}(t)=G_{12}(t)=(300t^3+345t^2+129t+16)/2$).

\emph{Step 2 (infinite domain).} Let $l\ge 27$ be odd. Then $l\ge27\ge12$ and $l-15\ge12$, so \eqref{eq:star} is valid at both $l$ and $l-15$, and both sides of the desired equality are given by branch polynomials: writing $l=60t+r$ with $r$ odd, the left side is $G_r(t)$ and the right side is $G_{r-15}(t)$ (with the wrap convention above), which are equal by Step 1.

\emph{Step 3 (finite base).} The remaining odd $l\in\{17,19,21,23,25\}$ are verified by exact enumeration of box partitions:
\[
g_6(17)=1=g_6(2),\quad g_6(19)=2=g_6(4),\quad g_6(21)=3=g_6(6),
\]
\[
g_6(23)=4=g_6(8),\quad g_6(25)=6=g_6(10).\qedhere
\]
\end{proof}

\begin{remark}[Machine-verified content of Theorem~\ref{thm:shift}]
Every step except the $30$ symbolic identities and the $5$ base values is a hand proof. The $30$ identities are identities of explicit polynomials of degree $\le4$ with denominators dividing $120$; they can be re-verified by hand or by any exact-arithmetic system from the branch data printed by the certificate. The theorem makes no asymptotic or probabilistic claim.
\end{remark}

\section{Central dominance and the equivalence theorem}\label{sec:cubic}

\subsection{A sound positivity test for cubics}
\begin{lemma}[Minimum test]\label{lem:cubic}
Let $C(t)\in\mathbb Z[t]$ have degree $\le 3$, and let $T_0\in\mathbb Z$. If $\deg C=3$ require its leading coefficient to be positive; if $\deg C=2$ require its leading coefficient to be positive; if $\deg C\le1$ require its slope to be nonnegative (a constant $C\ge0$ is accepted). Then $C(t)\ge0$ for all integers $t\ge T_0$ if and only if $C\ge0$ at the finite candidate set
\[
\{T_0\}\ \cup\ \bigl\{\lfloor\rho\rfloor,\lceil\rho\rceil\;:\;\rho \text{ a real root of } C',\ \rho> T_0\bigr\}\cap\{t\in\mathbb Z:t\ge T_0\}.
\]
(If the leading-coefficient condition fails for $\deg\ge1$, then $C(t)<0$ for arbitrarily large $t$ and the certificate is rejected.)
\end{lemma}
\begin{proof}
Under the leading-coefficient conditions $C(t)\to+\infty$ (or $C$ is nondecreasing), so a minimum over integers $t\ge T_0$ exists. $C'$ has degree $\le2$; between consecutive real roots of $C'$ the function $C$ is monotone on the reals, hence on the integers of each interval of monotonicity $\cap[T_0,\infty)$ the minimum is attained at an endpoint. The endpoints are $T_0$ or the integers adjacent to a root $\rho>T_0$. The rejection clause is immediate from the leading term.
\end{proof}

\begin{remark}\label{rem:cubic-history}
A weaker test (value, slope, and leading coefficient at a single point $t=T_0$) suffices for quadratics but is \emph{false} for cubics, and the tail certificates of this paper do contain genuine cubics: $2400t^3+420t^2-30t-1$ is one of the difference polynomials arising at $(6,8)$, $t\ge2$. All certificates in this paper use Lemma~\ref{lem:cubic}; the quadratic-only variant is not relied upon anywhere.
\end{remark}

\subsection{Tail certificates}
Fix $b\in\{8,12,16\}$, $h=b^2/8$, and let $l\ge122$. Repeating the argument of Lemma~\ref{lem:star} for the margin $M_j(l)$ at offsets $j\in\{-1,\dots,h-2\}$: $M_j(l)$ is an explicit $\{-1,0,+1\}$-linear combination of differences $D(y)=Q_6$-type values at arguments $3l-h+j$, $2l-i-h+j$, $l-i'-j-h+j$; for $l\ge122$ all surviving arguments are $\ge0$ and all size-$\ge3$ subsets vanish (the size-$3$ arguments are at most $-h+j-\sigma\le -2-6<0$; the larger subsets are smaller still; the pair arguments are at least $l-h-12\ge 122-32-12>0$). Hence for each residue $r\bmod 60$ and each offset $j\in[-1,h-2]$ there is an explicit polynomial $M_{j,r}(t)$ with
\[
M_j(r+60t)=M_{j,r}(t)\qquad(t\ge2),
\]
and the comparison polynomial
\[
C_{j,r}(t)\;:=\;M_{j,r}(t)-M_{h-1,r}(t)
\]
has degree $\le 3$: the branch polynomials are constructed symbolically and the certificate reports the observed maximum degree $3$ across all $3480$ comparison polynomials (the a-priori bound from Lemma~\ref{lem:qp-classical} is $4$; the drop is empirical structure, and Lemma~\ref{lem:cubic} is stated for degree $\le3$ so the certificate consumes exactly what is produced). The certificate verifies, for every family, every $r\bmod60$, and every $j\in\{-1,\dots,h-2\}$, that $C_{j,r}$ passes Lemma~\ref{lem:cubic} at $T_0=2$:
\[
\text{(6,8): } 480/480,\qquad \text{(6,12): } 1080/1080,\qquad \text{(6,16): } 1920/1920 \text{ checks pass.}
\]
For $l=r+60t\ge122$ this gives $M_j(l)\ge M_{h-1}(l)$ for all $j\le h-2$, hence
\begin{equation}\label{eq:tailred}
\textstyle\bigl(\min_j M_j(l)\ge 0\bigr)\iff\bigl(M_{h-1}(l)\ge 0\bigr)\iff\bigl(g_6(l)\ge\Delta_b\bigr),\qquad l\ge122 .
\end{equation}

\subsection{Finite core and the proof of Theorem~\ref{thm:equivalence}}
For each family, exact dynamic programming over box partitions verifies the equivalence ``$F^{(6,b)}_l$ unimodal $\Leftrightarrow$ all $M_j(l)\ge0$ $\Leftrightarrow$ $M_{h-1}(l)\ge0$'' for every admissible $l\le121$:
(6,8): $l=3$--$121$; (6,12): $l=6$--$121$; (6,16): $l=11$--$121$; zero mismatches. (The upper limit $121$ closes the seam: the per-class tail certificates are stated for $t\ge2$, i.e.\ every $l=r+60t\ge120$, and the core covers $l\le121$, so no integer is missed; the overlap at $l=120,121$ is redundant but harmless.) (The middle equivalence is the nontrivial content: it fails, e.g., at $(6,24)$, $l=46$.) Combining with \eqref{eq:tailred}:
\begin{proof}[Proof of Theorem~\ref{thm:equivalence}]
For $l\ge122$, \eqref{eq:tailred} plus Lemma~\ref{lem:window}. For admissible $l\le121$, the exact finite verification plus Lemma~\ref{lem:window}.
\end{proof}

\subsection{Clearance of the even class}\label{sec:clearance}
For the classification we need the barrier $g_6(l)\ge\Delta_b$ on the even classes. The values of $g_6$ on even $l$ (Table~\ref{tab:g6}) are strictly increasing through $l\le 40$; for the infinite tail we certify, class by class, that $G_r(t)\ge\Delta_b$ for all even residues $r$ and $t\ge2$ using Lemma~\ref{lem:cubic} applied to $G_r(t)-\Delta_b$ ($30$ polynomials per threshold), and by exact DP for even $l\le120$. Result: $g_6(l)\ge\Delta_8$ for all even $l\ge2$; $g_6(l)\ge\Delta_{12}$ for all even $l\ge6$; $g_6(l)\ge\Delta_{16}$ for all even $l\ge12$; all clearances pass ($90/90$ tail checks plus finite bases).

\begin{table}[h]
\centering
\begin{tabular}{c|ccccccccccccccc}
$l$&2&4&6&8&10&12&14&15&17&19&21&23&25&27&\dots\\ \hline
$g_6(l)$&1&2&3&4&6&8&10&1&1&2&3&4&6&8&\dots
\end{tabular}
\caption{Exact central gaps: even classes (increasing) and the odd class translated by $15$ ($g_6(27)=g_6(12)=8$, etc.). Odd $l\in\{1,\dots,13\}$ all have $g_6(l)=0$; $g_6(15)=1=g_6(0)$ by convention.}
\label{tab:g6}
\end{table}

\section{Proof of the classification}\label{sec:classification}
\begin{proof}[Proof of Theorem~\ref{thm:classification}]
By Theorem~\ref{thm:equivalence}, for each $b\in\{8,12,16\}$ and every admissible $l$, failure is equivalent to $g_6(l)<\Delta_b$.

\emph{Even $l$:} by \S\ref{sec:clearance}, $g_6(l)\ge\Delta_b$ for every even admissible $l$ (clearance starts at even $l\ge2$, $l\ge6$, $l\ge12$ respectively, and every admissible even $l$ is at or beyond the start, since admissibility forces $l\ge3$, $l\ge6$, $l\ge11$). No even $l$ fails.

\emph{Odd $l\ge17$:} by Theorem~\ref{thm:shift}, $g_6(l)=g_6(l-15)$ with $l-15$ even, so failure is equivalent to $g_6(l-15)<\Delta_b$. By Table~\ref{tab:g6} and clearance, this means $l-15<l_{0e}$, i.e.\ $l<l_{0e}+15$.
For $b=8$: $l_{0e}=0$, so no odd $l\ge17$ fails; the admissible odd $l<17$ are $3,5,\dots,15$, whose gaps are $g_6=0$ for $l\le13$ and $g_6(15)=1\ge1$; hence $E(6,8)=\{3,5,7,9,11,13\}$.
For $b=12$: $l_{0e}=6$; odd failures with $l\ge17$ are $17,19$ (translates $2,4$ with gaps $1,2<3$), plus odd admissible $l\in\{7,\dots,15\}$ with gaps $0,0,0,0,1<3$: $E(6,12)=\{7,9,\dots,19\}$.
For $b=16$: $l_{0e}=12$; odd failures with $l\ge17$ are $17\le l\le 25$ (translates $2,\dots,10$ with gaps $1,2,3,4,6<7$), plus odd admissible $l\in\{11,13,15\}$ with gaps $0,0,1<7$: $E(6,16)=\{11,13,\dots,25\}$.

Sharpness: $l_0-1=l_{0e}+13$ is odd and fails. For $b=12,16$ (where $l_{0e}\ge6$), its Shift-Law translate is $l_0-1-15=l_{0e}-2$, an \emph{even} parameter below $l_{0e}$, so by minimality of $l_{0e}$ its gap is below the barrier: concretely $g_6(4)=2<3$ for $b=12$ and $g_6(10)=6<7$ for $b=16$. For $b=8$ one has $l_{0e}=0$ and the claim is direct: $l_0-1=13$ with $g_6(13)=0<1$. Every admissible $l\ge l_0$ passes: even $l$ by \S\ref{sec:clearance}, odd $l$ because $l-15\ge l_0-15=l_{0e}-1$ forces the even translate $l-15\ge l_{0e}$ (parity), which clears by clearance. The identity $l_0=l_{0e}+14$ reads off Table~\ref{tab:g6}: $l_{0e}=0,6,12$ for $b=8,12,16$.
\end{proof}

\begin{remark}
The classification shows the exceptional sets are \emph{parity-pure}: contiguous odd segments at the admissibility boundary. This is the $k=6$ instance of the \emph{conjectured} phenomenon that for even $k$ the slow parity class is $kl\equiv2\pmod4$; the broader picture across $26$ families is exact-computation evidence only, not theorem \cite{M08}.
\end{remark}

\section{The constant 15 and parity}\label{sec:why15}
\subsection{What the proof explains}
The branch data make the rigidity visible: for every odd residue $r$, the branch polynomial $G_r(t)$ of the odd class is literally the even-class branch $G_{r-15}(t)$. The constant $15=\binom 62=5+4+3+2+1+0$ is the staircase sum of the box height. Among odd shifts $c\in\{1,3,\dots,45\}$, only $c=15$ satisfies $g_6(l)=g_6(l-c)$ for every odd $l$ in the test window $[c+17,129]$ \cite{M09}; this uniqueness is \emph{numerically} established and is not a theorem.

\subsection{Parity}
Across the even-$k$ families examined in \cite{M08}, the central gap of $\ggbrack{l+k}{k}$ is anomalously small precisely when $kl\equiv2\pmod4$; this parity pattern is \emph{observed and conjectural for $k>6$}, and only the $k=6$ instance is used (and proved) in this paper. Here $k=6$ and odd $l$ give $6l\equiv2$ or $6l\equiv 6\pmod{12}$, and the parity of the center $kl/2$ interacts with the half-integer arguments of the inclusion--exclusion branches. The Shift Law is a parity \emph{transport}: it maps the slow class ($l$ odd) onto the fast class ($l-15$ even) in one direction only. The inverse fails: $g_6(24)=29\ne0=g_6(9)$.

\subsection{Why naive bijections fail}\label{sec:naive}
\begin{enumerate}
\item \textbf{Staircase subtraction.} For $l\ge6$, subtracting $\delta=(5,4,3,2,1,0)$ bijects \emph{strict} partitions with exactly $6$ positive parts and largest part at most $l$ onto partitions into exactly $6$ positive parts with largest part at most $l-5$; the inverse adds $\delta$, positivity and the width bound are preserved on both sides, and the total weight drops by exactly $15$. The construction nevertheless fails for our purpose twice over: an arbitrary (weak) six-part partition is not mapped to a partition at all when consecutive parts repeat, so the exact-$6$-row summand of $a_l(3l)$ is only partially touched (partitions with fewer than $6$ parts are unpaired as well, and no sign-reversing pairing of the boundary layer was found); and the induced width shift is $l\mapsto l-5$, not the desired $l\mapsto l-15$.
\item \textbf{Complementation} $\lambda\mapsto(6l-|\lambda|)$ fixes the central weight and does not transport the width at all.
\item \textbf{2-core / rim-hook transports} meet runner boundaries imposed by the box: a $15$-rim hook is not available on every relevant diagram, and the core effect depends on its position.
\item Any bijective proof must be \emph{directional} (\S\ref{sec:why15}): no global period-$15$ symmetry exists (the even-parity comparison fails, e.g.\ $g_6(24)=29\ne0=g_6(9)$), so an argument that averages over a period-$15$ group action cannot apply.
\end{enumerate}

As observed in external review (2026-11-13), the \emph{representation-theoretic} mechanism is classical: $g_6(l)$ is the dimension of the degree-$l$ component of the invariant ring of binary sextics, which is generated by invariants of degrees $2,4,6,10$ together with one odd generator $E$ of degree $15$ subject to a single quadratic relation $E^2=P(S,T,H,I)$ \cite{GY}; the Hilbert series is
\[
\sum_{l\ge0} g_6(l)\,t^l=\frac{1+t^{15}}{(1-t^2)(1-t^4)(1-t^6)(1-t^{10})},
\]
and multiplication by $E$ identifies the even degree-$(l-15)$ component with the odd degree-$l$ component. The same generating function follows from Appendix A.2(1)--(2) of \cite{KKM}. This is precisely the residue-free explanation of the constant $15$; our quasipolynomial proof is an elementary self-contained route to the same identity.

\begin{conjecture}[Partition-level mechanism; open]\label{conj:mechanism}
There is a \emph{manifestly positive combinatorial} explanation of the identity $G_r\equiv G_{r-15}$ --- e.g.\ a signed transport pairing the boundary layers of $a_l(3l)-a_l(3l-1)$, or an explicit expansion of the relation $E^2=P(S,T,H,I)$ into partition moves --- that does not proceed by residue-class enumeration. The classical invariant-theoretic identification above is representation-theoretic, not a partition bijection; see \texttt{proofs/MECHANISM\_ATTEMPT.md} in the accompanying package for the strongest partial attempt found so far.
\end{conjecture}

\section{Related work and novelty position}\label{sec:related}
\textbf{Stanley--Zanello \cite{SZ}.} Conjecture 1.1 of \cite{SZ} concerns $f(k,m,b)=\ggbrack{m}{k}-q^{k(m-b)/2+b-2k+2}\ggbrack{b}{k-2}$ and asserts eventual unimodality for $k\ge5$ uniformly in $b$ with one explicit exception. Our family $\ggbrack{l+6}{6}-q^{3l-h}\ggbrack{b}{b/2}$ is a different comparison (the subtractand is central, of the same rectangle shape language as \cite{CJ}); the two families overlap only partially, and neither theorem implies the other. \cite{SZ} proves the conjecture for $k\le5$, where the classification of exceptional parameters is exact; our Theorem~\ref{thm:classification} is an exact classification inside a single family for $k=6$, $b\le16$.

\textbf{Koutschan--Uncu--Wong \cite{KUW}.} Their proof of the $k=6,7$ cases of \cite{SZ} uses the same broad technology --- coefficient sequences of rational generating functions are quasipolynomial, and positivity is decided per residue class, there by cylindrical algebraic decomposition. Our certificates are the same method specialized to fixed offsets with an explicit minimum test (Lemma~\ref{lem:cubic}); we claim no methodological novelty.

\textbf{Reiner--Stanton \cite{RS}.} Their Theorem 5 gives two-step propagation in the width for fixed subtractands, which reduces each parity tail of our classification to finitely many consecutive base cases once unimodality (not the gap equivalence) is the target. Our equivalence theorem is a different statement: it identifies the failure mode with a single central inequality \emph{and its exact domain of validity}, including the certified failure at $(6,24)$.

\textbf{Krattenthaler--Kronholm--Marsh \cite{KKM}.} Theorems 2.2/2.3 of \cite{KKM} give rational generating functions for fixed-center-offset coefficient sequences of $\ggbrack{N+m}{m}$, which independently justify the quasipolynomial structure used in \S\S\ref{sec:qp}--\ref{sec:cubic} for the $m=6$ case; our floor-bookkeeping certificate is self-contained. More strongly, Appendix A.2(1)--(2) of \cite{KKM} yields, for the $m=6$ central-offset sequence, a generating function equivalent to $(1+t^{15})\big/\prod(1-t^{d})$, $d\in\{2,4,6,10\}$, which implies Theorem~\ref{thm:shift} directly; this connection was made explicit by external review (2026-11-13) and is the reason the novelty statement in \S\ref{sec:intro} is scoped to the equivalence and classification results in \S\ref{sec:setup}.

\textbf{Classical invariant theory \cite{GY}.} The invariant ring of binary sextics is generated by invariants of degrees $2,4,6,10,15$ with a single relation of degree $30$; its Hilbert series is $(1+t^{15})/((1-t^2)(1-t^4)(1-t^6)(1-t^{10}))$, and since $g_6(l)$ is the dimension of the degree-$l$ invariant component (equivalently, the middle adjacent-coefficient difference of $\ggbrack{l+6}{6}$), the Shift Law for odd $l\ge15$ is immediate. Multiplication by the degree-$15$ generator is the mechanism. This classical identification, brought out by external review (2026-11-13), supersedes the earlier ``no mechanism known'' phrasing.

\textbf{Strict unimodality.} \cite{PP,Bounds,Dhand} lower-bound central gaps of a single Gaussian polynomial (gaps growing with the parameters); our $g_6$ is exactly such a gap, and the Shift Law says the odd-parity subsequence of this gap function is a translate of the even one --- a statement not of the strict-unimodality form.

\textbf{Companion results.} The window-reduction lemma in the shape used here first appeared in the companion paper \cite{M06} ($k=8$, $b=12$: $\ggbrack{l+8}{8}-q^{4l-18}\ggbrack{12}{6}$ unimodal iff $l\ge6$), whose infinite tail certificate uses period $840$; \cite{M07} certifies $(8,16)$, $(10,16)$, and the $(6,12)$ base block with thresholds agreeing with Theorem~\ref{thm:classification}. This paper is self-contained and no result here depends on \cite{M06,M07}.

\section{Limitations and open questions}\label{sec:limits}
\begin{enumerate}
\item The equivalence theorem holds, as proved, exactly on $\{(6,8),(6,12),(6,16)\}\times\{$admissible $l\}$. It is false for $b=24$ (Proposition~\ref{prop:ce}); the correct general formulation should include a curvature condition comparing the ascent profiles of the two rectangles. Extending the certificates to $(6,20)$ or $(6,24)$ with a window lemma is mechanical but not done.
\item For $k>6$ no analogue of the Shift Law is proved; at $k=8,10$ the corresponding shift candidates fail numerically \cite{M08}, and the general statement (if any) should involve $\binom k2$ and the slow parity class $kl\equiv2\pmod4$.
\item Uniqueness of the shift $15$ is verified numerically (\S\ref{sec:why15}), not proved.
\item The partition-level mechanism problem (Conjecture~\ref{conj:mechanism}) is open. The representation-theoretic explanation --- the invariant ring of binary sextics and multiplication by its degree-$15$ generator --- is classical (\S\ref{sec:why15}, \cite{GY,KKM}); what is missing is a manifestly positive combinatorial transport at the partition level.
\item The growth of $\Delta_b$ ($1,3,7,24,127,\dots$) has no closed form. It is at most exponential since $\Delta_b\le\binom{b}{b/2}\le 2^b$; the data are consistent with growth whose base approaches $2$, unproved.
\end{enumerate}

\section{Evidence classes and reproducibility}\label{sec:status}
\begin{center}\footnotesize
\begin{tabular}{p{0.46\textwidth}lp{0.2\textwidth}}
\toprule
Item & Class & Script \\
\midrule
Lemma~\ref{lem:window} (window reduction) & hand proof & --- \\
Lemma~\ref{lem:qp-classical}, \ref{lem:qp-cert}, \ref{lem:star}, \ref{lem:cubic} & hand proof & --- \\
$60/60$ recurrence identities + base $n\le119$ & certificate & \texttt{verify\_shift.py S1} \\
$30/30$ branch identities \eqref{eq:branch} & certificate & \texttt{verify\_shift.py S3} \\
Base cases $l=17$--$25$; Table~\ref{tab:g6} & exact finite computation & \texttt{verify\_shift.py S4} \\
Tail $C_{j,r}\ge0$: $480{+}1080{+}1920$ checks & certificate & \texttt{verify\_central.py} \\
Finite core $l\le121$ equivalence & exact finite computation & \texttt{verify\_core.py} \\
Even-class clearance $90/90$ + bases & certificate & \texttt{verify\_clearance.py} \\
Proposition~\ref{prop:ce} & exact finite computation & \texttt{counterexample\_6\_\allowbreak 24.py} \\
Uniqueness of shift $15$ & numerical evidence only & \texttt{shift\_uniqueness.py} \\
\bottomrule
\end{tabular}
\end{center}
All scripts use exact integer or exact rational arithmetic (Python~3, sympy); no floating point appears in any certificate. Corruption tests (a deliberately perturbed branch polynomial, a perturbed base value, a perturbed threshold) are detected with nonzero exit codes; see \texttt{REPRODUCIBILITY.md} of the accompanying package for commands, software versions, runtimes, and complete PASS/FAIL output.

\section{Appendix: sample certificate data}\label{app:cert}
\textbf{Branch polynomials (Shift Law).} With $l=60t+r$:
\begin{align*}
G_{27}(t)&=G_{12}(t)=\tfrac{1}{2}\left(300t^3+345t^2+129t+16\right),\\
G_{17}(t)&=G_{2}(t)=\tfrac{1}{2}\left(300t^3+195t^2+39t+2\right).
\end{align*}
At $t=0$: $G_{27}(0)=8=g_6(27)$ and $G_{12}(0)=8=g_6(12)$; $G_{17}(0)=1=g_6(17)=g_6(2)$. The remaining $28$ identities are produced in identical form by \texttt{verify\_shift.py}.

\textbf{A cubic in the central certificates.} At $(6,8)$, class $r=40$, offset $j=3$, the comparison polynomial is $C(t)=2400t^3+420t^2-30t-1$, which has a negative value at $t=0$ but is certified nonnegative for integers $t\ge2$ by Lemma~\ref{lem:cubic}: the critical points of $C$ are the roots of $240t^2+28t-1$, namely $(-7\pm\sqrt{109})/120$, both below $2$, so the candidate set reduces to $\{t=2\}$ (exact check; no derivative root exceeds $2$). This is the polynomial that invalidated the quadratic-only test in an earlier version of the certificate.

\textbf{Counterexample data.} $[q^{132}]F^{(6,24)}_{46}=165876$, $[q^{133}]F^{(6,24)}_{46}=165790$ (difference $-86$ at offset $j=66=h-6$), while $g_6(46)=\Delta_{24}=127$. At $l=48$: $[q^{142}]F^{(6,24)}_{48}=215262$, $[q^{143}]F^{(6,24)}_{48}=215251$ (difference $-11$ at offset $j=70=h-2$), while $g_6(48)-\Delta_{24}=14$.

\section*{Acknowledgments and disclosure}
The preparation of this manuscript and its verification package used AI assistance
(an OpenClaw agent working under the author's direction); the human author is
accountable for all statements. External review correspondence of 2026-11-13 is
documented in the accompanying package (\texttt{HERMES\_REVIEW.md}). The paper has
not undergone human peer review; machine-verified evidence classes are recorded in
\S\ref{sec:status}.

\begin{thebibliography}{99}
\bibitem{sylvester} J.\ J.\ Sylvester, \emph{Proof of the hitherto undemonstrated fundamental theorem of invariants}, Philos.\ Mag.\ (5) \textbf{5} (1878), 178--188; reprinted in Coll.\ Math.\ Papers \textbf{3}, Chelsea, 117--126. (Contains the classical unimodality of Gaussian polynomials; citation form as in \cite{SZ,KUW}.)
\bibitem{stanley} R.\ P.\ Stanley, \emph{Enumerative Combinatorics}, vol.\ 2, Cambridge Univ.\ Press, 1999.
\bibitem{ohara} K.\ M.\ O'Hara, \emph{Unimodality of Gaussian coefficients: a constructive proof}, J.\ Combin.\ Theory Ser.\ A \textbf{53} (1990), 29--52. DOI 10.1016/0097-3165(90)90018-r (Crossref-verified).
\bibitem{SZ} R.\ P.\ Stanley and F.\ Zanello, \emph{A generalization of a 1998 unimodality conjecture of Reiner and Stanton}, J.\ Combinatorics \textbf{11} (2020), 111--126; preprint arXiv:1711.10033 (v2, 2018). (Crossref-verified.)
\bibitem{KUW} C.\ Koutschan, A.\ K.\ Uncu and E.\ Wong, \emph{A Unified Approach to Unimodality of Gaussian Polynomials}, ISSAC 2023, ACM, DOI 10.1145/3597066.3597113; expanded version arXiv:2302.04067 (theorem numbers cited here refer to the arXiv version). (Crossref-verified.)
\bibitem{RS} V.\ Reiner and D.\ Stanton, \emph{Unimodality of differences of specialized Schur functions}, J.\ Algebraic Combin.\ \textbf{7} (1998), 91--107. DOI 10.1023/A:1008623312887 (Crossref-verified).
\bibitem{PP} I.\ Pak and G.\ Panova, \emph{Strict unimodality of $q$-binomial coefficients}, C.\ R.\ Acad.\ Sci.\ Paris \textbf{351} (2013), 415--418; corrected version arXiv:1306.5085.
\bibitem{Bounds} I.\ Pak and G.\ Panova, \emph{Bounds on Kronecker and $q$-binomial coefficients}, arXiv:1410.7087.
\bibitem{Dhand} V.\ Dhand, \emph{A combinatorial proof of strict unimodality for $q$-binomial coefficients}, Discrete Math.\ \textbf{335} (2014), 20--24; arXiv:1402.1199. (Crossref-verified.)
\bibitem{CJ} W.\ Y.\ C.\ Chen and I.\ D.\ D.\ Jia, \emph{Semi-invariants of binary forms pertaining to a unimodality theorem of Reiner and Stanton}, Internat.\ J.\ Math.\ \textbf{32} (2021), 2140003. DOI 10.1142/S0129167X21400036; arXiv:2011.01467. (Crossref-verified.)
\bibitem{KKM} C.\ Krattenthaler, B.\ Kronholm and P.\ Marsh, \emph{An orthogonal view of Gau\ss{}ian polynomials}, arXiv:2510.14124 (v2, 2026).
\bibitem{GY} J.\ H.\ Grace and A.\ Young, \emph{The Algebra of Invariants}, Cambridge University Press, 1903. (Author/title/year verified against library catalog records (Open Library entry for John Hilton Grace and Alfred Young, 1903; the digitized scan is held as University of Michigan Historical Math Collection record ABS4236.0001.001, not directly inspected here). The sextic generator description and Hilbert series used here are the classical formula; no specific page locator is claimed, and the series was re-checked against our exact coefficient data for $0\le l\le400$ and against \cite{KKM} Appendix A.2.)
\bibitem{KR} G.\ E.\ Andrews, \emph{The Theory of Partitions}, Cambridge Univ.\ Press, 1998.
\bibitem{M06} G.\ W.\ Hu, \emph{Unimodality of $\ggbrack{l+8}{8}-q^{4l-18}\ggbrack{12}{6}$ for all $l\ge6$}, unpublished companion preprint, 2026; no stable public identifier located. Not relied upon: this paper is self-contained.
\bibitem{M07} G.\ W.\ Hu, \emph{Certified sharp thresholds for centered Gaussian subtraction families}, unpublished companion preprint, 2026; no stable public identifier located. Not relied upon: this paper is self-contained.
\bibitem{M08} G.\ W.\ Hu, \emph{Central gaps, parity, and exceptional sets for two-rectangle subtractions}, unpublished companion preprint, 2026; cited only as the record of exploratory computations, which carry no theorem status here.
\bibitem{M09} G.\ W.\ Hu, \emph{The Shift Law for the central gap of $\ggbrack{l+6}{6}$}, unpublished companion preprint, 2026; cited only as the record of the numerical uniqueness scan (\S\ref{sec:why15}), which is verified-not-proved.
\end{thebibliography}

\end{document}
